% AC-03 — Artefato Central: Inventário de variáveis com fonte, tipo e cobertura
% Gerado manualmente em 2026-03-22

\chapter*{Apêndice B: Inventário de Variáveis do GSS}
\addcontentsline{toc}{chapter}{Apêndice B: Inventário de Variáveis do GSS}
\label{apendice:variaveis}

Este apêndice documenta todas as variáveis utilizadas no cálculo do
\textbf{Global Safety Score (GSS)} para os 262 distritos do Paraguai.
Para cada variável são indicados: a fonte primária de coleta, o tipo de
dado (direto, interpolado ou proxy), a cobertura geográfica efetiva e o
grau de confiança da estimativa.

\medskip
\noindent\textbf{Legenda de tipos:}
\begin{itemize}[nosep, leftmargin=1.5em]
  \item \textbf{Direto} — dado observado ou medido diretamente na fonte.
  \item \textbf{Interpolado} — calculado a partir de dados de unidades regionais
        adjacentes ou superiores.
  \item \textbf{Proxy} — variável substituta utilizada pela ausência de dado
        primário disponível.
\end{itemize}

\medskip
\noindent\textbf{Legenda de confiança:}
\begin{itemize}[nosep, leftmargin=1.5em]
  \item \textbf{A} — dado primário oficial verificado.
  \item \textbf{B} — estimativa interpolada de regionais oficiais.
  \item \textbf{C} — proxy inferido por analogia.
\end{itemize}

% ─────────────────────────────────────────────────────────────────────────────
\section*{Dimensão A — Ameaças Estratégicas (peso: 25\%)}
\addcontentsline{toc}{section}{Dimensão A — Ameaças Estratégicas}
\label{apendice:variaveis:dimA}

Esta dimensão avalia a exposição de cada distrito a ameaças de natureza
geopolítica e militar: proximidade de alvos estratégicos, porosidade de
fronteiras e vetores de instabilidade regional provenientes do crime
organizado transnacional.

\medskip
{\footnotesize\setlength{\tabcolsep}{3pt}\renewcommand{\arraystretch}{1.25}
\begin{tabular}{@{}>{\raggedright}p{2.8cm} >{\raggedright}p{2.6cm} >{\centering}p{1.3cm} >{\raggedright}p{2.0cm} >{\centering}p{1.0cm}@{}}
\toprule
\textbf{Variável} & \textbf{Fonte} & \textbf{Tipo} & \textbf{Cobertura} & \textbf{Conf.} \\
\midrule
Ausência de alvos militares estratégicos &
  Análise geopolítica pública &
  Proxy &
  Nacional uniforme &
  B \\
\midrule
Distância de fronteiras ativas &
  Cálculo geoespacial (DGEEC shapefile CNPV 2022) &
  Direto &
  262 distritos &
  A \\
\midrule
Risco de spillover do narcotráfico &
  Ministerio Público PY / InSight Crime &
  Proxy &
  Dept. fronteiriços &
  B \\
\midrule
Distância de Itaipu / Yacyretá (alvo industrial) &
  Cálculo geoespacial (coordenadas públicas) &
  Direto &
  262 distritos &
  A \\
\bottomrule
\end{tabular}%
}

% ─────────────────────────────────────────────────────────────────────────────
\section*{Dimensão B — Risco Social (peso: 20\%)}
\addcontentsline{toc}{section}{Dimensão B — Risco Social}
\label{apendice:variaveis:dimB}

Mensura criminalidade, desigualdade econômica, vulnerabilidade sistêmica
prisional e densidade de ameaças sociais decorrentes de grandes aglomerados
urbanos. As taxas de homicídio e os dados da EPHC constituem o núcleo
empírico mais robusto desta dimensão.

\medskip
{\footnotesize\setlength{\tabcolsep}{3pt}\renewcommand{\arraystretch}{1.25}
\begin{tabular}{@{}>{\raggedright}p{2.8cm} >{\raggedright}p{2.6cm} >{\centering}p{1.3cm} >{\raggedright}p{2.0cm} >{\centering}p{1.0cm}@{}}
\toprule
\textbf{Variável} & \textbf{Fonte} & \textbf{Tipo} & \textbf{Cobertura} & \textbf{Conf.} \\
\midrule
Taxa de homicídios (por 100k hab) &
  Ministerio Público PY (série 2019--2024) &
  Direto &
  18 departamentos &
  A \\
\midrule
Índice de Gini &
  INE EPHC 2023 &
  Direto &
  18 departamentos &
  A \\
\midrule
Taxa de pobreza total (\%) &
  INE EPHC 2023 &
  Direto &
  18 departamentos &
  A \\
\midrule
Taxa de pobreza extrema (\%) &
  INE EPHC 2023 &
  Direto &
  18 departamentos &
  A \\
\midrule
\% de presos provisórios &
  MJT / Ministerio de Justicia 2023 &
  Direto &
  18 departamentos &
  A \\
\midrule
Taxa de superlotação prisional &
  MJT 2023 &
  Direto &
  18 departamentos &
  A \\
\midrule
Presença de crime organizado transnacional &
  InSight Crime / Ministerio Público PY &
  Proxy &
  18 departamentos &
  B \\
\midrule
Distância de cidades $>$500k hab &
  Cálculo geoespacial &
  Direto &
  262 distritos &
  A \\
\bottomrule
\end{tabular}%
}

% ─────────────────────────────────────────────────────────────────────────────
\section*{Dimensão C — Riscos Naturais e Clima (peso: 15\%)}
\addcontentsline{toc}{section}{Dimensão C — Riscos Naturais e Clima}
\label{apendice:variaveis:dimC}

Cobre os principais vetores de risco ambiental: regime pluviométrico,
irradiância, temperatura, inundações recorrentes, seca e sismicidade.
Os dados NASA POWER oferecem cobertura distrital completa; os demais
variáveis dependem de interpolação ou análise hidrológica regional.

\medskip
{\footnotesize\setlength{\tabcolsep}{3pt}\renewcommand{\arraystretch}{1.25}
\begin{tabular}{@{}>{\raggedright}p{2.8cm} >{\raggedright}p{2.6cm} >{\centering}p{1.3cm} >{\raggedright}p{2.0cm} >{\centering}p{1.0cm}@{}}
\toprule
\textbf{Variável} & \textbf{Fonte} & \textbf{Tipo} & \textbf{Cobertura} & \textbf{Conf.} \\
\midrule
Precipitação média anual (mm/ano) &
  NASA POWER (2001--2020) &
  Direto &
  262 distritos &
  A \\
\midrule
Irradiância solar global média (kWh/m²/dia) &
  NASA POWER (2001--2020) &
  Direto &
  262 distritos &
  A \\
\midrule
Temperatura máxima absoluta histórica &
  SEAM / DINAC série histórica &
  Direto &
  18 capitais dept. &
  A \\
\midrule
Risco de inundação recorrente &
  SEAM / BID análise hidrológica &
  Proxy &
  18 departamentos &
  B \\
\midrule
Risco de seca severa &
  Precipitação $<$800 mm/ano como proxy &
  Proxy &
  262 distritos &
  B \\
\midrule
Risco sísmico &
  USGS Seismic Hazard Model SA &
  Direto &
  Nacional &
  A \\
\bottomrule
\end{tabular}%
}

% ─────────────────────────────────────────────────────────────────────────────
\section*{Dimensão D — Autossuficiência (peso: 20\%)}
\addcontentsline{toc}{section}{Dimensão D — Autossuficiência}
\label{apendice:variaveis:dimD}

Avalia a capacidade de um território sustentar produção alimentar, acesso
autônomo à água e geração energética local. O Censo 2022 do INE fornece
cobertura distrital para água e saneamento; as variáveis de solo provêm
majoritariamente de grades globais do SoilGrids (ISRIC), com resolução de
250 m.

\medskip
{\footnotesize\setlength{\tabcolsep}{3pt}\renewcommand{\arraystretch}{1.25}
\begin{tabular}{@{}>{\raggedright}p{2.8cm} >{\raggedright}p{2.6cm} >{\centering}p{1.3cm} >{\raggedright}p{2.0cm} >{\centering}p{1.0cm}@{}}
\toprule
\textbf{Variável} & \textbf{Fonte} & \textbf{Tipo} & \textbf{Cobertura} & \textbf{Conf.} \\
\midrule
Aptidão agrícola do solo (\% classes I--III) &
  SoilGrids ISRIC + SEAM &
  Direto / Interpolado &
  18 departamentos &
  B \\
\midrule
pH médio do solo &
  SoilGrids ISRIC 250 m &
  Direto &
  262 distritos (grid) &
  A \\
\midrule
Fertilidade natural do solo &
  SoilGrids: matéria orgânica + nutrientes &
  Direto &
  262 distritos (grid) &
  A \\
\midrule
Aquífero principal (tipo e profundidade) &
  SENASA / MADES-DGPCRH &
  Direto &
  18 departamentos &
  A \\
\midrule
Cobertura do Aquífero Guarani &
  ISARM / OAS (cartografia PY) &
  Direto &
  262 distritos (grid) &
  A \\
\midrule
Acesso a água potável (\% domicílios) &
  INE Censo 2022 &
  Direto &
  262 distritos &
  A \\
\midrule
Acesso a saneamento (\% domicílios) &
  INE Censo 2022 &
  Direto &
  262 distritos &
  A \\
\midrule
Irradiação solar para FV (kWh/m²/dia) &
  NASA POWER (2001--2020) &
  Direto &
  262 distritos &
  A \\
\midrule
Preço médio da terra rural (USD/ha) &
  INDERT / Mercado imobiliário 2024 &
  Proxy &
  18 departamentos &
  B \\
\bottomrule
\end{tabular}%
}

% ─────────────────────────────────────────────────────────────────────────────
\section*{Dimensão E — Ambiente Institucional (peso: 20\%)}
\addcontentsline{toc}{section}{Dimensão E — Ambiente Institucional}
\label{apendice:variaveis:dimE}

Captura qualidade de vida, infraestrutura digital, regime tributário e
estabilidade do marco legal para o estrangeiro residente. O IDH
departamental e os indicadores de educação e saúde do INE constituem as
fontes primárias; a cobertura 4G e a segurança jurídica são estimadas por
dados de operadoras e análise documental.

\medskip
{\footnotesize\setlength{\tabcolsep}{3pt}\renewcommand{\arraystretch}{1.25}
\begin{tabular}{@{}>{\raggedright}p{2.8cm} >{\raggedright}p{2.6cm} >{\centering}p{1.3cm} >{\raggedright}p{2.0cm} >{\centering}p{1.0cm}@{}}
\toprule
\textbf{Variável} & \textbf{Fonte} & \textbf{Tipo} & \textbf{Cobertura} & \textbf{Conf.} \\
\midrule
IDH departamental &
  PNUD / Global Data Lab 2020 &
  Direto &
  18 departamentos &
  A \\
\midrule
Esperança de vida ao nascer &
  INE Censo 2022 / PNUD &
  Direto &
  18 departamentos &
  A \\
\midrule
Anos médios de escolaridade &
  INE Censo 2022 &
  Direto &
  18 departamentos &
  A \\
\midrule
\% domicílios com internet &
  INE EPHC 2023 &
  Direto &
  18 departamentos &
  A \\
\midrule
Cobertura 4G Tigo (\% pop) &
  CONATEL dados de operadoras 2023 &
  Direto &
  18 departamentos &
  B \\
\midrule
Tributação efetiva (IRACIS flat 10\%) &
  SET / MH Paraguai &
  Direto &
  Nacional uniforme &
  A \\
\midrule
Segurança jurídica da propriedade &
  INDERT / Registro Público &
  Proxy &
  18 departamentos &
  B \\
\midrule
Estabilidade constitucional &
  Freedom House / análise histórica &
  Proxy &
  Nacional &
  B \\
\bottomrule
\end{tabular}%
}

% ─────────────────────────────────────────────────────────────────────────────
\section*{Lacunas Identificadas e Prioridades de Coleta}
\addcontentsline{toc}{section}{Lacunas Identificadas e Prioridades de Coleta}
\label{apendice:variaveis:lacunas}

A cobertura atual do dataset apresenta quatro lacunas estruturais que
limitam a granularidade distrital do GSS. As prioridades de coleta estão
ordenadas por impacto na precisão do score:

\begin{enumerate}[nosep, leftmargin=1.8em]
  \item \textbf{Criminalidade por distrito (262):} o Ministerio Público PY
        divulga séries por departamento (18 unidades). A desagregação
        distrital exigiria acesso à base de ocorrências do sistema SIGHPU
        — \emph{prioridade alta}.
  \item \textbf{Pobreza e renda por distrito:} a EPHC utiliza amostra
        representativa somente ao nível departamental. A expansão da
        amostra ou o uso de dados sintéticos do Censo 2022 seria necessária
        para cobertura distrital — \emph{prioridade alta}.
  \item \textbf{Preços de terra por distrito:} o INDERT e o mercado
        imobiliário dispõem de dados apenas por departamento ou por oferta
        pontual. Uma pesquisa de campo sistemática (amostragem de anúncios
        georreferenciados) poderia cobrir os 262 distritos —
        \emph{prioridade média}.
  \item \textbf{IDH por distrito:} o PNUD não calcula IDH abaixo do nível
        departamental no Paraguai. Um proxy via Índice de Condições de Vida
        (ICV) a partir do Censo 2022 é tecnicamente viável e recomendado
        para versões futuras do modelo — \emph{prioridade média}.
\end{enumerate}

% ─────────────────────────────────────────────────────────────────────────────
\section*{Cobertura Atual do Dataset}
\addcontentsline{toc}{section}{Cobertura Atual do Dataset}
\label{apendice:variaveis:cobertura}

O quadro abaixo sintetiza a distribuição de qualidade dos dados por
dimensão. Das 35 variáveis que compõem o modelo GSS, 23 (66\%) possuem
dado primário oficial verificado (confiança A), e 12 (34\%) dependem de
interpolação ou proxy (confiança B). Nenhuma variável foi classificada
como confiança C na versão atual.

\medskip
{\small
\renewcommand{\arraystretch}{1.3}
\begin{tabularx}{\linewidth}{>{\raggedright}X
                              >{\centering}p{2.0cm}
                              >{\centering}p{1.8cm}
                              >{\centering}p{1.8cm}
                              >{\centering\arraybackslash}p{1.4cm}}
\toprule
\textbf{Dimensão} & \textbf{Conf.\ A} & \textbf{Conf.\ B} & \textbf{Conf.\ C} & \textbf{Total} \\
\midrule
A — Ameaças Estratégicas (25\%) & 2 & 2 & 0 & 4 \\
B — Risco Social (20\%)         & 6 & 2 & 0 & 8 \\
C — Clima / Naturais (15\%)     & 4 & 2 & 0 & 6 \\
D — Autossuficiência (20\%)     & 7 & 2 & 0 & 9 \\
E — Institucional (20\%)        & 4 & 4 & 0 & 8 \\
\midrule
\textbf{Total}                  & \textbf{23} & \textbf{12} & \textbf{0} & \textbf{35} \\
\bottomrule
\end{tabularx}
}

\medskip
\noindent A proporção de 66\% de dados com confiança A posiciona o modelo
GSS-Paraguai como um dos mais bem fundamentados na literatura de
relocação estratégica aplicada. As dimensões B (Risco Social) e D
(Autossuficiência) concentram o maior número de variáveis de alta
confiança, refletindo a qualidade dos dados do Censo 2022 do INE e das
séries do Ministerio Público. A dimensão E (Institucional) apresenta a
maior proporção relativa de estimativas (50\% confiança B), o que indica
espaço para refinamento via dados primários adicionais de CONATEL e do
Registro Público.
